\chapter{吊钩端感知系统设计与标定}\label{ux540aux94a9ux7aefux611fux77e5ux7cfbux7edfux6784ux5efaux4e0eux6807ux5b9a}

本章针对吊钩端近场危险区域感知的工程实现，提出多传感器系统的总体架构、硬件集成与标定配准方案，并建立可供后续算法模块直接调用的数据接口与坐标系规范。由于后续的点云融合、定位建图与雷视关联均依赖传感器外参稳定性与跨模态时空一致性，本章的核心目标是将“可部署”的传感器平台转化为“可计算”的统一观测空间。

具体而言，\secref{ux7cfbux7edfux603bux4f53ux67b6ux6784ux4e0eux786cux4ef6ux9009ux578bux5e03ux7f6e}阐述吊钩端硬件选型与布置原则，分析在吊钩端特殊场景下如何实现视场覆盖与结构可靠性；\secref{ux6570ux636eux94feux8defux4e0eux63a5ux53e3ux89c4ux8303}规范通信链路、时间同步与延迟补偿策略，以减少链路抖动与时延带来的融合误差；\secref{sec:multi_sensor_coord_calib}进一步标定相机内参、双雷达外参与雷达--相机外参，并建立统一的坐标体系与TF接口规范，确保点云与图像的几何一致性，为第3章的融合点云定位与检测算法提供稳定输入。

\section{硬件系统架构与布置}\label{ux7cfbux7edfux603bux4f53ux67b6ux6784ux4e0eux786cux4ef6ux9009ux578bux5e03ux7f6e}

\subsection{硬件设备选型}\label{sec:hardware_selection}

吊钩端感知系统需在有限安装空间内覆盖吊钩周围近场危险区域，并在粉尘、强光/弱光、遮挡与振动等工况下保持稳定输出。为兼顾几何测距精度、环境鲁棒性与语义识别能力，本文采用双 LiDAR 与工业相机的异构配置：

双激光雷达（Livox Mid-360）作为近场空间几何感知的主传感器。LiDAR 具备 360° 水平覆盖能力，对光照变化不敏感，可在尘土等不利条件下保持稳定的测距输出。采用双雷达布置可通过不同俯仰角与基线配置实现盲区互补，并提升关键方向的点云密度与冗余度。

工业相机（海康威视 MV-CS016-10UC，镜头 MVL-HF0628M-6MPE，1/1.8''、600万像素）用于获取吊装构件与周围环境的图像观测。视觉信息与 LiDAR 的几何测距结果形成互补，支撑后续的吊装构件识别、精细化定位与环境理解，使系统既具备高精度的空间约束，也具备可解释的视觉证据。

系统的 IMU 信息来源于 Mid-360 内置 IMU，用于运动约束与点云去畸变，从而在吊钩端摆动与回转工况下维持位姿估计与点云时序对齐的稳定性。

系统的核心计算单元采用 NVIDIA Jetson Orin NX 嵌入式 AI 计算平台，负责所有传感器数据的采集、预处理与融合计算。该模块搭载 8 核 Arm Cortex-A78AE CPU 与 1024 核 Ampere 架构 GPU（含 32 个 Tensor Core），板载 16 GB 128 位 LPDDR5 内存，整机功耗在 15-25W 范围内，满足吊钩端电池供电场景下的续航要求。

为展示本文所采用的关键硬件设备实物，图\ref{fig:lidar_device_photo}--图\ref{fig:jetson_device_photo}分别给出了激光雷达、工业相机、工业镜头和边缘计算盒的实物图。

\begin{figure}[htbp]
  \centering
  \includegraphics[width=0.5\textwidth]{激光雷达-实物.png}
  \caption{激光雷达 Livox Mid-360}
  \label{fig:lidar_device_photo}
\end{figure}

\begin{figure}[htbp]
  \centering
  \includegraphics[width=0.5\textwidth]{工业相机-实物.png}
  \caption{工业相机 MV-CS016-10UC}
  \label{fig:camera_device_photo}
\end{figure}

\begin{figure}[htbp]
  \centering
  \includegraphics[width=0.5\textwidth]{工业镜头-实物.png}
  \caption{工业镜头 MVL-HF0628M-6MPE}
  \label{fig:lens_device_photo}
\end{figure}

\begin{figure}[htbp]
  \centering
  \includegraphics[width=0.5\textwidth]{jetson边缘盒.png}
  \caption{集成边缘计算盒}
  \label{fig:jetson_device_photo}
\end{figure}

为便于工程表达，表\ref{tab:sensor_spec}列出了本系统关键硬件组件及其在方法体系中的作用说明。

\begin{table}[htbp]
  \centering
  \caption{吊钩端感知系统关键硬件组件与作用}
  \begin{tabular*}{\textwidth}{@{\extracolsep{\fill}}>{\raggedright\arraybackslash}p{5cm}>{\raggedright\arraybackslash}p{2cm}>{\raggedright\arraybackslash}p{6.3cm}}
    \hline
    \textbf{组件} & \textbf{作用} & \textbf{关键说明}\\
    \hline
    LiDAR-1（Mid-360） & 近场几何感知 & 作为主测距源，提供近场点云，用于ROI建模、距离阈值判定与目标几何状态估计\\
    LiDAR-2（Mid-360） & 补盲与冗余 & 与LiDAR-1形成盲区互补，提高关键方向点云覆盖与稳定性，支撑融合点云生成\\
    工业相机（MV-CS016-10UC） & 语义识别 & 检测吊装构件并输出相关运动状态信息，用于雷视关联与后续避障逻辑\\
    镜头（MVL-HF0628M-6MPE） & 成像质量 & 满足作业距离下的分辨率需求，提升弱光/逆光条件下可用性\\
    IMU（内置） & 运动约束 & 为位姿估计提供短时高频运动信息，辅助点云去畸变与定位稳定性\\
    Jetson Orin NX & 边缘计算 & 16 GB 统一内存，搭载 GPU 加速感知管线，承载算法推理与计算任务\\
    \hline
  \end{tabular*}
  \label{tab:sensor_spec}
\end{table}

\subsection{传感设备安装布局设计}\label{ux5b89ux88c5ux5e03ux5c40ux4f18ux5316}

吊钩端传感器布局旨在实现近场空间全覆盖与跨模态数据关联。针对吊钩作业时吊钩块、吊具、钢丝绳及吊物本体造成的复杂遮挡，需要通过空间互补布置在 20m 近场范围内建立稳定可靠的感知覆盖。

双 LiDAR 布置遵循盲区互补、结构遮挡规避与振动冲击可控的原则。为直观说明单个激光雷达的感知边界，图\ref{fig:single_lidar_fov_placeholder}分别从竖直与水平方向展示了其视场角范围。单个 LiDAR 存在固有的视野盲区，无法满足吊钩端全方位监测需求。为此，本文采用两台 Mid-360 在盒体左右两侧侧向对称布置，安装高度与俯仰姿态保持一致，形成镜像观测关系。两台雷达均设置为$38^{\circ}$ 俯仰角，这一角度设计在避免激光互相干扰的前提下，显著增强了对吊钩下方近场区域的覆盖能力。

\begin{figure}[htbp]
  \centering
  \begin{subfigure}[t]{0.44\textwidth}
    \centering
    \includegraphics[height=4cm,keepaspectratio]{激光雷达角度示意图-水平方向.pdf}
    \caption{水平方向视场角}
    \label{fig:single_lidar_fov_horizontal}
  \end{subfigure}
  \hfill
  \begin{subfigure}[t]{0.54\textwidth}
    \centering
    \includegraphics[height=3.6cm,keepaspectratio]{激光雷达角度示意图-竖直方向.pdf}
    \caption{竖直方向视场角}
    \label{fig:single_lidar_fov_vertical}
  \end{subfigure}
  \caption{单个激光雷达视场角示意}
  \label{fig:single_lidar_fov_placeholder}
\end{figure}

工业相机的主要观测目标是吊装构件本身。相机采用垂直向下的安装方式，对吊钩下方作业面进行俯视拍摄，重点覆盖吊装构件区域。这一布置方案能够最大程度减少吊索对镜头的遮挡，在相对开阔的区域获取清晰的吊装构件图像。图\ref{fig:hook_platform_sensor_range}展示了双 LiDAR 与工业相机的整体探测范围及其空间重叠关系。从覆盖分布看，两台 LiDAR 的探测区域共同指向吊钩下方近场工作空间，形成几何感知的主覆盖区；工业相机则聚焦于核心作业面，与双 LiDAR 的高密度观测区域形成空间重叠。

\begin{figure}[htbp]
  \centering
  \includegraphics[width=0.9\textwidth]{吊钩感知平台传感器探测范围示意图.pdf}
  \caption{吊钩感知平台传感器探测范围示意}
  \label{fig:hook_platform_sensor_range}
\end{figure}

这种布局的核心价值在于建立点云与像素的跨模态关联。通过雷视外参标定建立的坐标变换关系，系统能够将点云中的三维几何特征与图像中的像素信息精确关联，从而实现对吊装构件空间状态（位置、姿态、运动趋势）的联合估计。这一关联关系为后续避障逻辑判断提供了关键基础：几何点云提供了精确的距离约束与空间边界，视觉图像则提供了场景理解与状态验证的可解释证据。

\subsection{吊钩端雷视融合一体化感知平台设计}\label{ux540aux94a9ux7aefux4e00ux4f53ux5316ux76d2ux4f53ux8bbeux8ba1}

为满足施工现场工程化部署需求，吊钩端感知系统采用一体化盒体进行集成，将传感器、电池组与计算模块集中布置，保障各传感器稳定固定安装以完成预期工作任务。在传感器布置方式与安装位置选择上，本文将传感设备集成安装于感知平台，该平台直接吊装于塔吊自身吊钩下方，而非零散安装于吊钩横梁，主要基于以下考虑：其一，塔吊吊钩横梁形式因设备型号与厂家而异，不存在统一的安装接口与平面，难以设计通用的传感器固定方案，且吊钩上方存在吊索、钢丝绳及起重臂等结构遮挡，若将传感器布置于上方将导致吊物正下方区域出现严重观测盲区；其二，下方安装使传感器随吊物同步运动，保证对吊物及近场环境的持续覆盖；其三，传感器与感知平台保持固定的相对位姿关系，是多传感器外参标定可执行的基本前提。图\ref{fig:sw_platform_overall}展示了感知平台的安装位置与整体外观。

\begin{figure}[htbp]
  \centering
  \includegraphics[width=0.88\textwidth]{传感平台整体外观示意图.pdf}
  \caption{吊钩下方雷视融合一体化集成感知平台安装位置与整体外观}
  \label{fig:sw_platform_overall}
\end{figure}

基于上述设计约束，本文采用 SolidWorks 完成吊钩下方传感器平台的结构设计与装配验证。在结构实现上，吊钩端平台采用上盖与下盖的两体式盒体方案：下盖承担传感器与支架的主要安装基准，上盖用于防护与密封，并与下盖形成封闭腔体以容纳电源与计算单元。盒体内部集成电池、NVIDIA Jetson 边缘计算模块、走线与固定件等部件，所有供电与信号线缆均在盒体内完成布线，同时对外部保留必要的充电/调试接口。此外，盒体在设计时预留了多处固定螺栓孔洞，为现场安装及后期维护调整提供灵活的接口余量。为直观展示上盖/下盖分体结构及各部件装配关系，图\ref{fig:sw_platform_exploded}给出了主要部件装配关系示意。在此基础上，图\ref{fig:sw_platform_layout}进一步给出了内部关键部件的布局与线束走线方案，图中标注了电池、嵌入式计算平台及供电/通信模块的安装位置，该布局在满足散热与接口可达性的前提下，保证供电与通信链路组织清晰，为现场组装与检修预留操作空间。

\begin{figure}[htbp]
  \centering
  
\includegraphics[width=0.88\textwidth]{传感平台主要配件关系示意图.pdf}
  \caption{雷视融合一体化集成感知平台上盖、下盖与主要部件装配关系}
  \label{fig:sw_platform_exploded}
\end{figure}

\begin{figure}[htbp]
  \centering
  
\includegraphics[width=0.8\textwidth]{盒体内部布局与走线示意.pdf}
  \caption{雷视融合一体化集成感知平台内部关键部件布局}
  \label{fig:sw_platform_layout}
\end{figure}

在完成整体结构设计的基础上，进一步确定传感器与平台间的连接紧固方式。感知平台的安全性是需要优先满足的设计约束——平台不得干扰吊钩原有的荷载传递路径。吊物荷载由双吊钩直接承担，感知平台不参与任何承力，仅通过焊接于吊钩外侧螺母柱的矩形板件获得托举支撑。双吊钩之间通过六角型加长螺母柱及螺钉实现连接紧固。在传感器固定方面，双激光雷达通过M3螺丝紧固于下盖双侧的专用安装底座，工业相机通过M3螺丝固定于上盖侧端，上下盒体之间通过M4螺栓连接紧固。上述安装方式在保证传感器刚性与位姿稳定的前提下，确保各自视场覆盖且互不干扰，同时使感知平台不构成附加吊装安全风险。

\section{通信链路与接口设计}\label{ux6570ux636eux94feux8defux4e0eux63a5ux53e3ux89c4ux8303}

\subsection{通信链路设计与实现}\label{ux901aux4fe1ux94feux8defux642dux5efa}

吊钩端多传感器协同感知对通信可靠性与实时性均提出了严格要求。为此，本文设计如下通信架构：各传感器在物理层保持独立接入，由边缘计算单元统一完成数据采集与预处理，涵盖时间戳管理、数据缓存与消息发布等环节。系统总线将预处理后的数据解耦发布，实现底层感知链路与上层定位、检测及融合算法的分离。该设计在保障数据传输可靠性的同时，亦满足多传感器融合的实时性约束。图\ref{fig:hardware_link}展示了相应的硬件组成与通信关系，为后续的吞吐量约束与接口规范提供了结构基础。

\begin{figure}[htbp]
  \centering
  \includegraphics[width=0.92\textwidth]{吊钩端多传感器硬件与通信链路.pdf}
  \caption{吊钩端多传感器硬件与通信链路}
  \label{fig:hardware_link}
\end{figure}

链路性能瓶颈主要源于瞬时带宽波动、解码与内存拷贝开销及队列积压。本文通过吞吐量模型对通信负载施加约束：对任意话题 $i$，其平均吞吐量 $\bar{B}_i$ 可表示为
\begin{equation}
\bar{B}_i = f_i\,S_i
\end{equation}
其中，$f_i$ 为发布频率，$S_i$ 为单帧消息大小。系统总吞吐量 $\bar{B}=\sum_i\bar{B}_i$ 需保留峰值流量与协议开销裕量。该约束将链路可持续性转化为可量化判据，为频率配置与压缩策略提供统一依据。

跨模态链路采用“频率匹配”与“流量整形”的联合策略。在频率匹配层面，双路 LiDAR 的工作频率固定为 $10\,\mathrm{Hz}$——该取值在保证近场传感分辨率的同时，将链路带宽与边缘端计算负荷控制在可接受范围内。为降低跨模态时间对齐的复杂度、提升特征对应的稳定性，本文将相机视觉观测频率同步设定为 $10\,\mathrm{Hz}$，使两类传感器在统一的采样频率下工作。

以吞吐量公式对视觉链路进行量化说明：相机分辨率为 $1440\times1080$、RGB 三通道编码，单帧原始消息大小 $S_{\text{raw}}\approx4.7\,\text{MB}$，在 $f=10\,\mathrm{Hz}$ 下代入可得
\begin{equation}
  \bar{B}_{\text{raw}} = f\cdot S_{\text{raw}} = 10 \times 4.7\,\text{MB} \approx 47\,\text{MB/s}
\end{equation}

双路 LiDAR 点云链路的总吞吐量约为 $6\,\text{MB/s}$，双路 IMU 链路约为 $135\,\text{KB/s}$。上述传感器数据均以 $10\,\text{Hz}$ 固定频率持续流入边缘计算单元。处理管线各节点以流水线方式并行运行，同一时刻不同节点处理不同帧的数据，单个节点的消息队列为各帧提供缓冲去耦，管线在稳态下瞬时工作集约等于各传感器单帧数据量之和。中间各模块的输出以 ROS 话题形式在节点间传递；话题仅在有订阅者时产生实际传输开销，配合有界队列的流控策略，即使中间结果带宽较大，也不会造成总线拥塞或数据堆积。

网络侧采用单域二层汇聚链路，以减少不必要的转发路径与配置复杂度。各传感设备与边缘计算单元处于同一二层广播域内，数据经统一接入口汇聚至边缘侧，从而避免跨网段转发、桥接切换或多跳路径引入的附加时延。考虑到后续时间同步需要稳定的主从关系，网络中仅保留单一主时钟参考，并将链路状态、丢包统计与时间偏移估计纳入统一监测范围。该设计将数据接入、时间同步与状态监测纳入统一管理平面，从而降低时延抖动并提升长时运行稳定性。

ROS 话题通信采用发布/订阅机制，若生产者端速率高于消费者端处理能力，队列积压将直接转化为延迟增长，破坏多传感器融合的一致性。为此，本文对关键话题采用“小队列”与“可丢帧”策略，使系统尽可能运行在低延迟稳态。设订阅端平均处理时延为 $T_p$、队列长度为 $Q$，在持续拥塞且单帧处理时延近似恒定的条件下，队列引入的额外排队时延可近似表示为
\begin{equation}
T_q \approx (Q-1)\,T_p
\end{equation}

因此，对于点云、图像等数据量较大的消息，应依据链路功能分层配置队列参数。在当前工程实现中，点云与 IMU 聚合链路在融合节点侧采用较大的缓存队列，订阅队列与发布队列均配置为 $1000$，以提升突发负载下的稳定性。跨模态关联链路采用有限同步队列，图像与点云订阅队列配置为 $10$，近似时间同步器队列长度配置为 $20$，并将同步时间窗约束为 $0.05\,\mathrm{s}$，以抑制累积排队时延。

为统一接口语义并降低模块耦合，本文将关键数据流划分为原始采集接口与处理后输出接口两类。原始采集接口内容如表\ref{tab:ros_topics_min}所示：点云数据通过 livox\_ros\_driver2/CustomMsg 消息类型由激光雷达原生驱动发布；惯性测量单元（IMU）数据通过 sensor\_msgs/Imu 发布；视觉观测由海康相机原生驱动以 sensor\_msgs/Image 发布，并经压缩后以 sensor\_msgs/CompressedImage 分发；视觉标定信息通过 sensor\_msgs/CameraInfo 发布。处理后输出接口则主要包括融合点云、目标检测结果与定位位姿等结构化状态量，供下游规划与控制模块订阅。

\begin{table}[htbp]
  \centering
  \caption{吊钩端感知系统关键 ROS 接口内容}
  \label{tab:ros_topics_min}
  {\setlength{\tabcolsep}{3pt}
  \renewcommand{\arraystretch}{0.95}
  \begin{tabularx}{\textwidth}{L{2.1cm}L{4.7cm}L{2.1cm}>{\raggedright\arraybackslash}X}
    \toprule
    \textbf{数据类别} & \textbf{接口语义} & \textbf{输入频率} & \textbf{接口说明与作用} \\
    \midrule
    Lidar观测 & \mbox{livox\_ros\_driver2/CustomMsg} & $10\,\mathrm{Hz}\times2$ & 点云接口，提供空间几何观测；用于融合、ROI 建模与测距 \\
    IMU观测 & \mbox{sensor\_msgs/Imu} & $200\,\mathrm{Hz}\times2$ & 惯性接口，提供角速度与加速度；用于状态估计与点云去畸变 \\
    视觉观测 & \mbox{sensor\_msgs/CompressedImage} & $10\,\mathrm{Hz}$ & 图像接口，默认发布压缩相机帧；原始 RGB 按需启用，用于目标识别与雷视关联 \\
    视觉标定 & \mbox{sensor\_msgs/CameraInfo} & $10\,\mathrm{Hz}$ & 标定接口，提供内参与畸变参数；用于投影计算与几何一致性约束 \\
    \bottomrule
  \end{tabularx}
  }
\end{table}

\subsection{多传感器时间同步与延迟补偿}\label{ux65f6ux95f4ux540cux6b65ux4e0eux5ef6ux8fdfux8865ux507f}

多传感器协同感知以时间一致性为核心前提。受采样频率差异、缓存机制、链路传输与解码处理开销等因素影响，不同观测在时间轴上通常存在系统偏差与随机抖动。若直接以消息到达时刻进行融合，将导致雷视关联漂移与点云投影误差。针对该问题，本文以统一时基下的消息时间戳为参考，通过底层授时同步、上层时间对齐与延迟补偿实现时间一致性保障。

时间源采用分层设计思路：具备硬件授时能力的设备使用 PTP 时基，其余设备使用系统时钟时间戳。统一时基的关键在于以 PTP 作为主时间参考，在必要时将系统时钟与之对齐，再于融合层执行时间对齐与补偿，而非要求所有传感器采用同一硬件授时方式。

Livox 激光雷达支持 PTP、gPTP、GPS 等多种时间同步方式。就吊钩端部署条件而言，gPTP 主要面向 802.1AS 配置体系，在普通工业以太网链路中部署复杂度较高；GPS 授时需额外集成接收设备及对应天线，工程上不具备实施条件；硬件时间戳在长时采集场景下稳定性亦不足，存在跳变风险。综合考虑，本文采用 IEEE 1588-2008（PTPv2）协议，以 LinuxPTP 软件时间戳模式建立统一授时基准。

在 PTPv2 机制下，从设备与主时钟通过报文交互获得 $t_1,t_2,t_3,t_4$，并据此估计链路时延与主从时钟偏移。其中，$t_1$ 与 $t_4$ 分别表示主时钟发送 Sync 报文与主时钟接收 Delay\_Req 报文的时刻，$t_2$ 与 $t_3$ 分别表示从设备接收 Sync 报文与发送 Delay\_Req 报文的时刻。

\begin{figure}[htbp]
  \centering
  \includegraphics[width=0.88\textwidth]{{PTP 时间同步示意图}.pdf}
  \caption{PTP 时间同步示意图}
  \label{fig:ptp_timing}
\end{figure}

PTP时间同步原理如图\ref{fig:ptp_timing}所示，相应的传输路径延迟与时钟偏移计算式如下。
\begin{equation}
\text{Delay} = \frac{(t_4-t_1)-(t_3-t_2)}{2}
\end{equation}

\begin{equation}
\text{Offset} = (t_2-t_1)-\text{Delay} = \frac{(t_2-t_1)+(t_3-t_4)}{2}
\end{equation}

当 Offset 收敛并稳定在小幅波动区间时，可认为 LiDAR 点云时间戳已与主时钟保持一致。

同一二层网络中若出现多个主时钟源将导致主从关系振荡，引入额外的时间同步抖动。因此，本系统采用 PTP 软件授时，并在融合层执行时间对齐与补偿，从而保证跨传感器时间一致性。该方法实现的是统一时基下的软件一致性，并不等价于所有设备均具备硬件级同步能力。

需要强调的是，底层授时同步并不能替代融合层的时间对齐与补偿机制。即便 LiDAR 侧已完成 PTPv2 软件同步，ROS 融合侧仍需执行最近邻匹配、插值重采样与固定延迟补偿，以保证异构观测在同一物理时刻对齐。对工业相机链路而言，由于图像时间戳来自系统时钟采样，还需通过时间偏差估计与门控策略抑制跨模态对齐误差。

设 LiDAR 点云帧时间戳为 $t^{l}_k$，相机图像时间戳为 $t^{c}_m$，IMU 测量时间戳序列为 $t^{i}_n$。融合与关联通常需要将不同频率数据对齐到同一参考时刻 $t$：对给定 $t$，可采用最近邻对齐选择满足 $\lvert t^{s}-t\rvert$ 最小的传感器测量作为对齐样本（$s\in\{l,c,i\}$）；对于高频 IMU，当需要 $t$ 时刻的测量值时，可在相邻两帧 $(t^{i}_{n}, \mathbf{z}^{i}_{n})$ 与 $(t^{i}_{n+1}, \mathbf{z}^{i}_{n+1})$ 之间进行线性插值
\begin{equation}
\mathbf{z}^{i}(t) = \alpha\,\mathbf{z}^{i}_{n} + (1-\alpha)\,\mathbf{z}^{i}_{n+1},\quad
\alpha = \frac{t^{i}_{n+1}-t}{t^{i}_{n+1}-t^{i}_{n}}
\end{equation}

对于 LiDAR 点云，若驱动提供点时间（每点相对帧起始的时间偏移），则可在定位模块中结合 IMU 实施去畸变（deskew），以抑制吊钩摆动与回转导致的运动扭曲。

对于相机与 LiDAR 等异构传感器链路，端到端时延视为可测量、可更新的系统参数，通过日志时间戳统计获得估计 $\hat{\Delta}$，并在融合侧执行时间补偿：
\begin{equation}
\tilde{t} = t - \hat{\Delta}
\end{equation}
其中，$t$ 为原始时间戳，$\tilde{t}$ 为补偿后的有效时间戳，$\hat{\Delta}$ 为包含采集、传输、解码与调度处理等环节的时延估计值。

上述时间同步与延迟补偿逻辑如图\ref{fig:time_sync_timeline}所示。

\begin{figure}[htbp]
  \centering
  \begin{tikzpicture}[x=2.4cm,y=1.0cm,>=Latex,font=\scriptsize]
    % lane headers
    \node[draw, rounded corners=2pt, fill=gray!15, minimum width=1.95cm, minimum height=0.9cm, align=center] (h1) at (0,0) {主机时钟 /\\PTP master};
    \node[draw, rounded corners=2pt, fill=gray!15, minimum width=1.95cm, minimum height=0.9cm, align=center] (h2) at (1,0) {LiDAR /\\内置 IMU};
    \node[draw, rounded corners=2pt, fill=gray!15, minimum width=1.95cm, minimum height=0.9cm, align=center] (h3) at (2,0) {系统时钟};
    \node[draw, rounded corners=2pt, fill=gray!15, minimum width=1.95cm, minimum height=0.9cm, align=center] (h4) at (3,0) {USB 相机};
    \node[draw, rounded corners=2pt, fill=gray!15, minimum width=1.95cm, minimum height=0.9cm, align=center] (h5) at (4,0) {ROS\\消息时间戳};
    \node[draw, rounded corners=2pt, fill=gray!15, minimum width=1.95cm, minimum height=0.9cm, align=center] (h6) at (5,0) {融合层};

    % lifelines (extend lower to avoid clipping)
    \foreach \x in {0,1,2,3,4,5}{
      \draw[dashed, gray] (\x,-0.55) -- (\x,-13.30);
    }

    % --------------------------
    % Equal vertical spacing plan (step = 0.8)
    % --------------------------
    % PTP:   -1.6, -2.4, -3.2
    % notes: -4.0
    % system clock follow:-4.8 (note -5.5)
    % pub:   -6.4, -7.2, -8.0, -8.8
    % fusion: -9.6, -10.4, -11.2, -12.0
    % camera note: -12.75

    % background groups (resized to cover new y-range cleanly)
    \draw[rounded corners=4pt, dashed, gray!70, fill=gray!6] (-0.55,-0.95) rectangle (2.55,-5.90);
    \node[anchor=north, font=\scriptsize] at (1.0,-1.05) {底层授时与统一时基};

    \draw[rounded corners=4pt, dashed, gray!70, fill=gray!6] (3.45,-6.10) rectangle (5.55,-13.40);
    \node[anchor=north, font=\scriptsize] at (4.5,-6.20) {上层时间对齐与延迟补偿};

    % --------------------------
    % PTP sync (shift down by 0.4~0.8 to avoid overlap; keep equal spacing)
    % --------------------------
    \draw[->, thick] (0,-2.10) -- (1,-2.10) node[midway, above] {Sync/Follow\_Up};
    \draw[->, thick] (1,-2.90) -- (0,-2.90) node[midway, above] {Delay\_Req};
    \draw[->, thick, dashed] (0,-3.70) -- (1,-3.70) node[midway, above] {Delay\_Resp};

    % system-clock alignment (keep equal spacing with Delay_Resp)
    \draw[->, thick] (0,-4.50) -- (2,-4.50) node[midway, above] {可选：phc2sys 跟随};

    % device-side note (under system-clock alignment arrow)
    \node[draw, fill=white, align=center, inner sep=2pt] at (1.0,-5.25)
      {LiDAR 点云与内置 IMU\\对齐至设备侧 PTP 时基};

    % system clock note (slightly lower, avoid overlap with next arrows)
    % \node[draw, fill=white, align=center, inner sep=2pt] at (2,-6.35)
    %   {系统时钟对齐至 PTP 基准（可选，默认关闭）};

    % --------------------------
    % sensor publication (equal spacing)
    % --------------------------
    \draw[->, thick] (1,-7.20) -- (4,-7.20) node[midway, above] {点云帧 $t_k^l$};
    \draw[->, thick] (2,-8.00) -- (3,-8.00) node[midway, above] {IMU 序列 $t_n^i$};
    \draw[->, thick] (3,-8.80) -- (4,-8.80) node[midway, above] {系统时间戳};
    \draw[->, thick] (3,-9.60) -- (4,-9.60) node[midway, above] {图像帧 $t_m^c$};

    % --------------------------
    % fusion processing (equal spacing)
    % --------------------------
    \draw[->, thick] (4,-10.40) -- (5,-10.40) node[midway, above] {最近邻对齐到参考时刻 $t$};
    \draw[->, thick] (4,-11.20) -- (5,-11.20) node[midway, above] {IMU 插值 $\mathbf{z}^{i}(t)$};
    \draw[->, thick] (4,-12.00) -- (5,-12.00) node[midway, above] {LiDAR 点时间支持 deskew};
    \draw[->, thick] (4,-12.80) -- (5,-12.80) node[midway, above] {估计相机链路延迟 $\hat{\Delta}$};

    % gating/compensation boxes (shift right + slightly up to avoid overlaps)
    \node[draw, fill=white, align=center, inner sep=2pt, text width=2.35cm] at (4.85,-7.60)
      {时间窗门控\\$\lvert t_a-t_b\rvert\le \tau$};
    \node[draw, fill=white, align=center, inner sep=2pt, text width=2.35cm] at (4.85,-8.80)
      {时间补偿\\$\tilde{t}=t-\hat{\Delta}$};

    % camera note (moved lower + left a bit, narrower to avoid covering arrows)
    \node[draw, rounded corners=2pt, fill=white, text width=6cm, align=left] at (2.10,-12.9)
      {USB 相机不具备 PTP 硬同步能力；在本文中采用融合层时间对齐与延迟补偿维持软件时间一致性。};

  \end{tikzpicture}
  \caption{多传感器时间同步与延迟补偿时序示意}
  \label{fig:time_sync_timeline}
\end{figure}

\section{多传感器坐标体系与标定}\label{sec:multi_sensor_coord_calib}

为实现多传感器数据的统一表达与融合，本节首先建立系统的坐标系框架与 TF 约定，并在此基础上完成相机内参、双雷达相对外参以及雷达‑相机外参标定，为后续点云与图像的几何对齐提供基础。

\subsection{坐标系定义与变换发布约定}\label{ux5750ux6807ux7cfbux4e0e-tf-ux63a5ux53e3ux89c4ux8303}


坐标变换（Transform，简称 TF）用于描述不同参考系之间的刚体关系，是多传感器空间对齐的基础。与时间同步解决观测在时间轴上的一致性不同，坐标系定义解决的是不同传感器在同一几何框架下的表达问题。若点云、惯性测量单元（IMU）与视觉观测分别采用不一致的参考系，将导致距离阈值、感兴趣区域（ROI）边界与投影关系难以直接比较，进而影响定位精度与风险判定的可靠性。本文基于刚体变换理论建立统一的空间参考体系。

设三维空间中的两个坐标系 $\{a\}$ 与 $\{b\}$，它们之间的刚体变换可表示为 $SE(3)$ 群中的元素：
\begin{equation}
\mathbf{T}_b^a = \begin{bmatrix}
\mathbf{R}_b^a & \mathbf{t}_b^a \\
\mathbf{0}^{\mathsf{T}} & 1
\end{bmatrix} \in SE(3)
\end{equation}
其中，$\mathbf{R}_b^a \in SO(3)$ 为旋转矩阵，$\mathbf{t}_b^a \in \mathbb{R}^3$ 为平移向量。该变换将坐标系 $\{b\}$ 中的点 $\mathbf{p}^b$ 映射到坐标系 $\{a\}$ 中：
\begin{equation}
\tilde{\mathbf{p}}^a = \mathbf{T}_b^a \tilde{\mathbf{p}}^b = \mathbf{R}_b^a \mathbf{p}^b + \mathbf{t}_b^a
\end{equation}
其中，$\tilde{\mathbf{p}} = [\mathbf{p}^{\mathsf{T}}, 1]^{\mathsf{T}}$ 为齐次坐标表示。

基于上述数学框架，本文提出一种分层式空间参考体系，将动态定位与静态感知解耦。该体系由“动态定位主链”与“静态感知支链”构成：动态定位主链负责吊钩主体的连续位姿估计与运动跟踪，形成从全局地图到本体系的主体运动描述；静态感知支链则管理多传感器之间的固定几何关系，确保跨模态观测在统一参考系下的对齐一致性。这种分层设计使得运动估计与传感器外参管理相互独立，提高了系统的模块化程度与可维护性。

\begin{figure}[htbp]
  \centering
  \begin{minipage}{0.9\textwidth}
    \footnotesize
\begin{verbatim}
map                                    # 全局地图系
└── odom                               # 局部连续里程计系
  └── base_link                      # 主体参考系
    ├── body                       # 融合点云/IMU 参考帧
    │   └── body_leveled           # 重力对齐参考帧
    │       └── camera_leveled     # 视觉投影参考帧
    └── chassis_link               # 兼容结构支路
      └── imu_link
\end{verbatim}
  \end{minipage}
  \caption{系统实际 TF 树与坐标系关系}
  \label{fig:tf_tree_hook}
\end{figure}

在机器人操作系统（ROS）框架下，坐标变换通过 TF（Transform）库进行发布与查询。本文针对吊钩端近场感知的特定需求，制定了如图\ref{fig:tf_tree_hook}所示的坐标系定义与 TF 发布规范。该树状结构体现了坐标系的继承关系与变换链的可追溯性，保证了从全局地图到传感器坐标系的完整变换路径。

为明确各坐标系的语义定义与功能定位，表\ref{tab:coordinate_frames}详细列出了关键坐标系及其在感知系统中的角色。

\begin{table}[htbp]
  \centering
  \caption{关键坐标系定义与用途}
  \label{tab:coordinate_frames}
  {\setlength{\tabcolsep}{4pt}
  \renewcommand{\arraystretch}{1.1}
  \begin{tabularx}{\textwidth}{L{3.0cm}L{7.4cm}L{3.4cm}}
    \toprule
    \textbf{坐标系} & \textbf{描述} & \textbf{备注} \\
    \midrule
    map & 全局地图系，表示固定的世界坐标系 & 全局定位参考 \\
    odom & 局部连续里程计系，提供连续无跳变的局部运动估计 & 用于局部路径规划 \\
    base\_link & 吊钩主体参考系，通常与机械结构固联 & 主体运动基准 \\
    body & 融合点云与 IMU 的参考帧，与 base\_link 为恒等变换 & 多传感器融合基准 \\
    body\_leveled & 重力对齐后的参考帧，Z 轴与重力方向平行 & 视觉投影与地面估计 \\
    camera\_leveled & 相机投影参考帧，与 body\_leveled 保持固定变换 & 雷视关联与图像投影 \\
    chassis\_link & 兼容结构支路，用于扩展机械参考系 & 本系统中未实际使用 \\
    imu\_link & IMU 传感器坐标系 & 用于惯性测量对齐 \\
    \bottomrule
  \end{tabularx}
  }
\end{table}

传感器安装外参与中间参考帧之间的变换在系统运行期间视为静态刚体变换，这些变换通过第2.3节的标定过程获得。以融合参考系 $\{\text{body}\}$ 为基准，各传感器坐标系的相对位姿可形式化表示如下：
\begin{gather}
\mathbf{T}_{\text{body}}^{\text{lidar\_0}} = \mathbf{I}_{4\times 4} \\[5pt]
\mathbf{T}_{\text{body}}^{\text{lidar\_1}} = \hat{\mathbf{T}}_{\text{body}}^{\text{lidar\_1}} \\[5pt]
\mathbf{T}_{\text{body}}^{\text{imu\_0}} = \mathbf{I}_{4\times 4} \\[5pt]
\mathbf{T}_{\text{body}}^{\text{imu\_1}} = \hat{\mathbf{T}}_{\text{body}}^{\text{lidar\_1}}
\end{gather}
其中，$\mathbf{I}_{4\times 4}$ 为单位矩阵，$\mathbf{T}_{\text{body}}^{\text{lidar\_0}} = \mathbf{I}_{4\times 4}$ 表示基准雷达坐标系 LiDAR-0 坐标系与融合参考系重合，$\mathbf{T}_{\text{body}}^{\text{imu\_0}}$ 同理，但需要说明的是，Mid-360 内置 IMU 与 LiDAR 之间的刚体变换由 Livox SDK 在驱动层内部维护，应用层直接使用已对齐至 LiDAR 帧的 IMU 观测；$\hat{\mathbf{T}}_{\text{body}}^{\text{lidar\_1}}$ 为通过第\ref{sec:lidar_lidar_calib}节双雷达外参标定获得的刚体变换；$\mathbf{T}_{\text{body}}^{\text{imu\_1}}$ 采用与 $\mathbf{T}_{\text{body}}^{\text{lidar\_1}}$ 相同的外参矩阵，因为 IMU-1 与 LiDAR-1 刚性连接于同一安装底座，相对位姿固定不变。

为确保空间参考体系的一致性与可维护性，坐标系约定采用右手坐标系，其中 X 轴指向前方，Y 轴指向左侧，Z 轴指向上方；单位制统一规定长度单位采用国际单位制米（m），角度单位采用弧度（rad），确保数值计算的一致性；旋转表示采用四元数 $\mathbf{q}=[q_w, q_x, q_y, q_z]^{\mathsf{T}}$ 或旋转矩阵 $\mathbf{R}\in SO(3)$，避免欧拉角奇异性问题；坐标变换需通过 TF 库的统一接口进行查询。以链式变换为例，点从 body 系变换到 body\_leveled 系可表示为
\begin{equation}
\mathbf{p}^{\text{body\_leveled}} = \mathbf{T}_{\text{body}}^{\text{body\_leveled}}\,\mathbf{p}^{\text{body}}
\end{equation}

进一步变换到 camera\_leveled 系则表示为
\begin{equation}
\mathbf{p}^{\text{camera\_leveled}} = \mathbf{T}_{\text{body\_leveled}}^{\text{camera\_leveled}}\,\mathbf{p}^{\text{body\_leveled}}
\end{equation}

这种链式查询方式不仅降低了多次变换组合引入的数值误差，也简化了实现复杂度。在外参管理方面，本文规定各算法模块内部若需中间参考系，均应通过标准 TF 链组合获得，禁止维护独立于 TF 系统的私有外参，从而确保变换关系的单一来源。

本章提出的分层式空间参考体系为后续章节的点云融合定位、危险区域检测与雷视关联算法提供了统一、一致且可追溯的坐标框架，确保了算法模块间的空间可比性。

\subsection{相机内参标定}\label{ux76f8ux673aux5185ux53c2ux6807ux5b9a}

相机内参标定旨在获取焦距、主点与畸变参数，以确保图像去畸变与点云到图像投影的几何一致性。本文采用平面棋盘格标定板采集多组不同距离、不同姿态的标定图像，并基于张正友标定法求解相机内参矩阵 $\mathbf{K}$ 及径向/切向畸变系数。标定完成后，使用重投影误差评价标定质量。

相机内参标定以针孔成像模型描述三维点到像素点的投影关系。设空间点在相机坐标系下为 $\mathbf{P}_c=[X\;Y\;Z]^{\mathsf{T}}$，其在归一化成像平面上的坐标为 $\mathbf{p}_n=[x\;y]^{\mathsf{T}}$，则
\begin{equation}
\label{eq:camera_projection}
\begin{bmatrix}
x \\
y \\
1
\end{bmatrix}
= \frac{1}{Z}
\begin{bmatrix}
X \\
Y \\
Z
\end{bmatrix},\qquad
s\begin{bmatrix}
u \\
v \\
1
\end{bmatrix}
= \mathbf{K}
\begin{bmatrix}
x \\
y \\
1
\end{bmatrix}
\end{equation}
其中，$(u,v)$ 为像素坐标，$s$ 为尺度因子，相机内参矩阵 $\mathbf{K}$ 定义为
\begin{equation}
\label{eq:intrinsic_matrix}
\mathbf{K}=
\begin{bmatrix}
f_x & 0 & c_x \\
0 & f_y & c_y \\
0 & 0 & 1
\end{bmatrix}
\end{equation}

式中 $f_x,f_y$ 为像素单位的等效焦距，$(c_x,c_y)$ 为主点坐标。若以棋盘格平面作为世界坐标系 $\{w\}$ 的参考，则相机外参（标定时每张图像对应一组）可写为旋转矩阵 $\mathbf{R}$ 与平移向量 $\mathbf{t}$，有
\begin{equation}
\label{eq:extrinsic_matrix}
s\begin{bmatrix}
u \\
v \\
1
\end{bmatrix}
=\mathbf{K}\,[\mathbf{R}\;|\;\mathbf{t}]\begin{bmatrix}
X_w \\
Y_w \\
Z_w \\
1
\end{bmatrix}
\end{equation}

张正友标定法利用平面标定板在不同姿态下产生的多组单应性（Homography）约束，在无需精密三维标定场的条件下同时估计 $\mathbf{K}$ 与畸变参数，具有实现简单、精度可靠等工程优势。

实际镜头通常存在光学畸变，本文采用径向畸变与切向畸变的组合模型进行矫正。令 $r^2=x^2+y^2$，则畸变后的归一化坐标 $(x_d,y_d)$ 可表示为
\begin{equation}
\label{eq:distortion_model}
x_d = x\left(1+k_1 r^2+k_2 r^4+k_3 r^6\right)+2p_1xy+p_2\left(r^2+2x^2\right)
\end{equation}
\begin{equation}
\label{eq:distortion_model_y}
y_d = y\left(1+k_1 r^2+k_2 r^4+k_3 r^6\right)+p_1\left(r^2+2y^2\right)+2p_2xy
\end{equation}
其中，$k_1,k_2,k_3$ 为径向畸变系数，$p_1,p_2$ 为切向畸变系数。将 $(x_d,y_d)$ 代入式\eqref{eq:camera_projection}即可得到带畸变的像素坐标预测值。

在实现层面，本文相机内参标定使用 MATLAB（Computer Vision Toolbox）完成，采用棋盘格标定板进行数据采集与参数求解。选用平面棋盘格标定板并记录方格边长 $L$（mm），在固定焦距与曝光等关键参数的前提下采集了24张多姿态图像，覆盖不同距离、不同角度及不同画面位置，避免样本集中在图像中心附近或姿态变化过小。图\ref{fig:cam_calib_setup}展示了各种角度、远近的棋盘格图像采集情况。角点提取阶段在 MATLAB 中使用 Camera Calibrator 应用自动检测棋盘角点，核验每张图像角点是否完整、是否存在误检，并剔除模糊、强反光或角点缺失严重的样本。参数求解阶段在 MATLAB 中设置方格边长 $L$后计算 $\mathbf{K}$ 与畸变系数，并自动给出每张图像的外参。


\begin{figure}[htbp]
  \centering
  \includegraphics[width=0.88\textwidth]{相机内参标定的数据采集示意}
  \caption{相机内参标定数据采集}
  \label{fig:cam_calib_setup}
\end{figure}

标定质量评价采用重投影误差衡量标定精度。设第 $i$ 张图像上第 $j$ 个角点的观测像素为 $\mathbf{p}_{ij}$，模型预测像素为 $\hat{\mathbf{p}}_{ij}$，则单点重投影误差为 $\lVert \mathbf{p}_{ij}-\hat{\mathbf{p}}_{ij} \rVert_2$，总体均方根误差（RMSE）定义为
\begin{equation}
\label{eq:reproj_error}
e_{\mathrm{rms}}=\sqrt{\frac{1}{N}\sum_{i=1}^{M}\sum_{j=1}^{n_i}\left\lVert \mathbf{p}_{ij}-\hat{\mathbf{p}}_{ij} \right\rVert_2^2}
\end{equation}
其中，$M$ 为图像数量，$n_i$ 为第 $i$ 张图像可用角点数，$N=\sum_i n_i$ 为总角点数。MATLAB 通常同时给出每张图像的重投影误差分布，便于发现离群样本并迭代剔除。误差统计与可视化结果在文中预留，如图\ref{fig:cam_calib_reproj}所示。

\begin{figure}[H]
  \centering
  \includegraphics[width=0.88\textwidth]{相机内参标定的重投影误差评价示意}
  \caption{相机内参标定的重投影误差评价}
  \label{fig:cam_calib_reproj}
\end{figure}

最终标定输出包括内参矩阵 $\mathbf{K}$、畸变系数 $\{k_1,k_2,k_3,p_1,p_2\}$ 以及每张标定图像对应的外参 $\{\mathbf{R}_i,\mathbf{t}_i\}$。本文将 $\mathbf{K}$ 与畸变参数用于图像去畸变与“点云到图像”的投影计算，为后续雷视外参标定与跨模态关联提供统一的几何基础。

\subsection{双雷达外参标定}\label{sec:lidar_lidar_calib}

双雷达外参标定的目标是在两台 Mid-360 之间建立稳定、可复用的刚体变换，使两路观测在后续融合、定位与雷视关联中共享统一的空间语义。本文求解的外参需满足自动精配准的收敛需求，并直接作为运行时融合节点的静态配置参数。为此，本节将双雷达标定明确划分为粗配准初值构造与自动精配准优化两个阶段：前者负责建立正确的坐标约定并提供可收敛初值，后者负责在真实点云数据上进一步压缩几何残差。

记基准雷达坐标系为 $\{l_1\}$，目标雷达坐标系为 $\{l_2\}$，待求外参为
\begin{equation}
\label{eq:lidar_lidar_extrinsic}
\mathbf{T}_{l_2}^{l_1}=
\begin{bmatrix}
\mathbf{R} & \mathbf{t}\\
\mathbf{0}^{\mathsf{T}} & 1
\end{bmatrix}\in SE(3)
\end{equation}
其中，$\mathbf{R}\in SO(3)$，$\mathbf{t}\in \mathbb{R}^3$。对任一点 $\mathbf{p}^{l_2}$，其在 $\{l_1\}$ 中的表达为
\begin{equation}
\tilde{\mathbf{p}}^{l_1}=\mathbf{R}_{l_2}^{l_1}\,\mathbf{p}^{l_2}+\mathbf{t}_{l_2}^{l_1}
\end{equation}

这一表达不仅给出了目标雷达点到基准雷达坐标系的统一方式，更是粗配准阶段必须先固定下来的矩阵方向约定。本文统一以基准雷达 $\{l_1\}$ 作为双雷达公共参考系，后续自动精配准与运行时调用均沿用该定义，因此本节仅讨论外参的离线构造与验证，不再展开在线阶段的坐标链组织。

粗配准阶段的任务并非追求最终精度，而是依据盒体结构设计、安装姿态与孔位基准给出物理上合理的初始位姿，其核心目的在于为精配准阶段的外参求解提供可靠的初值——初值不当将导致优化过程陷入局部最优解，甚至无法正确收敛。依据吊钩端一体化结构的 SolidWorks 模型，可得到双雷达安装底座的相对位置和姿态关系，并据此构造初始变换矩阵 $\mathbf{T}_0$。该初值随后被写入自动标定工具的初始化文件，用作后续逐帧配准的起点。

自动精配准部分参考 Livox 开源的多雷达外参自动标定工具实现思路\footnote{Livox\ Automatic\ Calibration\ Tools: \url{https://github.com/Livox-SDK/Livox_automatic_calibration}}，其核心流程为“基准雷达先建图，目标雷达逐帧与地图配准并迭代外参，最后基于一致性筛选得到稳定外参”。该方法要求初值具备合理的空间位置与重叠方向，且对时间同步与运动畸变较为敏感。为此，数据采集时应保持缓慢平稳运动，避免急转弯，并确保双雷达帧序列在时间上尽可能对齐。

自动标定首先在固定参考系 $\{w\}$ 下用基准雷达的多帧点云构建地图 $\mathcal{M}$，然后将目标雷达的每帧点云 $\mathcal{P}_k^{(2)}$ 通过当前外参估计变换到基准雷达或世界坐标下，并与地图 $\mathcal{M}$ 进行配准，得到增量变换 $\Delta \mathbf{T}_k$ 用于修正外参。本文采用广义 ICP（Generalized ICP, GICP）进行点云配准，在点到点最近邻基础上引入局部协方差近似点到平面的约束，以提升稀疏或噪声场景下的稳定性，其优化形式为
\begin{equation}
\label{eq:gicp_cost}
\min_{\mathbf{T}}\sum_{i}\left(\mathbf{p}_i-\mathbf{T}\mathbf{q}_i\right)^{\mathsf{T}}\mathbf{C}_i^{-1}\left(\mathbf{p}_i-\mathbf{T}\mathbf{q}_i\right)
\end{equation}
其中，$\mathbf{q}_i$ 为源点云中的点，$\mathbf{p}_i$ 为地图中与其对应的匹配点（或局部子地图点），$\mathbf{C}_i$ 为由两侧局部几何估计得到的协方差矩阵。优化得到的配准变换用于更新外参估计，例如
\begin{equation}
\label{eq:extrinsic_update}
\mathbf{T}_{l_2}^{l_1}\leftarrow \Delta \mathbf{T}_k\,\mathbf{T}_{l_2}^{l_1}
\end{equation}

通过配准质量评估机制，引入 fitness score 阈值、最大对应距离与迭代次数上限作为判别准则，可抑制错误匹配导致的外参估计发散。当配准质量不满足预设阈值时，系统将跳过当前帧的外参更新或回退至初始猜测值，以避免不可靠观测对估计结果的影响。逐帧配准过程产生外参估计序列 $\{\mathbf{T}_k\}$，该序列可能因动态目标、遮挡、运动畸变及局部几何退化等因素引入离群值。借鉴 Livox 开源标定工具的实现思路，在平移与欧拉角参数空间对 $\{\mathbf{T}_k\}$ 进行稳健拟合与一致性检验，筛选配准质量较高的帧并剔除离群估计，从而得到稳定的外参估计 $\hat{\mathbf{T}}$。该方法利用“刚性安装下外参为常值”的物理约束，使标定结果在长时序数据中具备更高的稳定性与可复现性。

数据准备阶段，本文选择结构稳定且环境几何特征丰富的室内场景进行数据采集，确保两台雷达在时间上尽可能对齐，并以较慢速度移动以降低单帧运动畸变。本文采用 ROS 录制两台 LiDAR 的 rosbag 数据，在离线阶段将两路 LiDAR 数据按固定时间片导出为 .pcd 帧，同时确保两路点云在帧序号与时间戳层面一一对应，保持统一的读取顺序与目录组织方式，以便后续自动标定工具的处理。数据整理与帧对齐的示意如图\ref{fig:lidar_calib_dataset}所示。

\begin{figure}[htbp]
  \centering
  \begin{minipage}{0.88\textwidth}
    \footnotesize
\begin{verbatim}
data
│
├── Base_LiDAR_Frames                # 基准激光雷达帧序列
│   ├── 100000.pcd                    # PCD点云文件，按序号命名
│   ├── 100001.pcd
│   ├── ...
│
├── Target_LiDAR_Frames               # 目标激光雷达帧序列
│   ├── 100000.pcd                    # 与基准雷达帧对应的点云文件
│   ├── 100001.pcd
│   ├── ...
│
├── parameters.txt                     # 配置文件
├── Init_Matrix.txt                    # 初始外参估计值
├── H_LiDAR_Map_data                   # 地图输出/缓存目录
│   ├── H_LiDAR_Map.pcd                # mapping生成的激光雷达地图点云
│
├── T_Matrix.txt                        # mapping生成的每帧4x4位姿矩阵
└── calib_data.txt                      # calibration生成的每帧标定日志
\end{verbatim}
  \end{minipage}
  \caption{双雷达外参自动标定数据准备}
  \label{fig:lidar_calib_dataset}
\end{figure}

粗配准外参 $\mathbf{T}_0$ 可由吊钩端盒体的结构设计模型和安装位姿计算得到，并在工程实现中写入初始化配置文件，作为逐帧配准更新的初始猜测。该初值不必达到最终标定精度，但必须保证两路点云在同一场景下具有足够重叠，以提高 GICP 收敛到正确极值的概率。需要指出的是，粗配准并不是“直接把 CAD 里的镜像关系抄成矩阵”，因为外参旋转必须满足 $\mathbf{R}\in SO(3)$，即只能由正规旋转组成，而不能把结构上的镜像关系误写成反射变换。因此，在给出初值矩阵前，需要先明确其旋转项的几何来源。

本文依据激光雷达间的刚体空间旋转与平移关系构造粗配准初值。以基准雷达 $\{l_1\}$ 为参考，根据 SolidWorks 模型中的设计值获取两台雷达的安装面法向、设备中心点及主视场方向。由于两台 Mid-360 在盒体左右两侧呈侧向镜像布置，且安装高度与俯仰姿态基本一致，粗配准旋转需满足两个几何约束：其一，两台雷达的横向语义相反，即目标雷达坐标系中的左右方向需翻转到与基准雷达一致；其二，两台雷达主视轴在盒体纵向剖面内存在固定夹角，需施加俯仰补偿将目标雷达的前向轴转换到基准雷达参考系下。若结构设计说明将该俯仰关系拆解为多个绕同一轴的连续角度，则应在参数层面合并为等效夹角 $\theta$ 后写入旋转矩阵。按照本文的坐标定义，先对目标雷达坐标系施加绕 $z$ 轴的 $180^{\circ}$ 旋转以完成镜像安装带来的轴向翻转，再依据 CAD 中测得的主视轴夹角 $\theta$ 施加绕 $y$ 轴的补偿旋转，从而得到粗配准旋转初值
\begin{equation}
\label{eq:lidar_init_rotation}
\mathbf{R}_0=\mathbf{R}_y(\theta)\,\mathbf{R}_z(\pi)=
\begin{bmatrix}
-\cos\theta & 0 & \sin\theta \\
0 & -1 & 0 \\
\sin\theta & 0 & \cos\theta
\end{bmatrix}
\end{equation}
其中，$\mathbf{R}_z(\pi)$ 对应镜像安装下的方位翻转，$\mathbf{R}_y(\theta)$ 对应两主视轴在 $x$-$z$ 剖面内的夹角补偿。本系统根据结构模型测得 $\theta\approx 76^{\circ}$，因此式\eqref{eq:lidar_init_rotation}中的 $-\cos\theta$ 与 $\sin\theta$ 分别对应后续矩阵中的 $-0.2419$ 与 $0.9703$。也就是说，初值矩阵中第二行的 $-1$ 并非经验设定，而是由双雷达左右对称安装导致的轴向翻转；而 $x$-$z$ 平面上的非零项则反映了两雷达主视轴之间的固定开角。

旋转初值确定后，平移项由两台雷达中心点在盒体坐标系中的相对位置给出，并统一到基准雷达参考系下。若 SolidWorks 中读出的基线偏置最初按目标雷达的局部轴定义，则在装配齐次矩阵前应先利用当前旋转初值完成参考系转换，记为
\begin{equation}
\label{eq:lidar_init_translation}
\mathbf{t}_0=\mathbf{R}_0\,\mathbf{t}_{\mathrm{local}}=
\begin{bmatrix}
\Delta x & \Delta y & \Delta z
\end{bmatrix}^{\mathsf{T}}
\end{equation}
其中，$\mathbf{t}_{\mathrm{local}}$ 表示目标雷达局部坐标系下定义的中心点偏置，$\Delta x,\Delta y,\Delta z$ 分别表示转换后在基准雷达参考系中的前后、左右和上下偏置。结合当前盒体的左右对称布置，$\Delta y$ 近似为 0，而 $\Delta x$ 与 $\Delta z$ 由 SolidWorks 模型中两底座中心点的几何距离直接读取。依据 SolidWorks 设计模型，可以得到理想状态的双雷达底座中心点间相对位置，如图\ref{fig:lidar_calib_estimate}所示，进而写出完整的粗配准初值矩阵如式\eqref{eq:lidar_init_guess}所示。

\begin{figure}[htbp]
  \centering
  \includegraphics[width=0.88\textwidth]{初始矩阵的位置估计}
  \caption{初始矩阵的位置估计}
  \label{fig:lidar_calib_estimate}
\end{figure}

\begin{equation}
\label{eq:lidar_init_guess}
\mathbf{T}_0=
\begin{bmatrix}
-\cos\theta & 0 & \sin\theta & \Delta x \\
0 & -1 & 0 & \Delta y \\
\sin\theta & 0 & \cos\theta & \Delta z \\
0 & 0 & 0 & 1
\end{bmatrix}
=
\begin{bmatrix}
-0.2419 & 0 & 0.9703 & 0.1395 \\
0 & -1 & 0 & 0 \\
0.9703 & 0 & 0.2419 & -0.1090 \\
0 & 0 & 0 & 1
\end{bmatrix}
\end{equation}

随后，以基准雷达帧序列构建地图 $\mathcal{M}$ 作为后续配准参考；在外参迭代阶段，将目标雷达每帧点云在当前外参下变换后与地图或局部子地图进行 GICP 配准，得到增量修正 $\Delta \mathbf{T}_k$ 并更新外参，同时记录每帧配准质量指标以支撑一致性筛选。标定过程中前后点云叠加效果如图\ref{fig:lidar_calib_overlay}所示。

\begin{figure}[htbp]
  \centering
  \includegraphics[width=0.88\textwidth]{双雷达自动标定的点云配准效果示意.png}
  \caption{双雷达自动标定的点云配准效果示意}
  \label{fig:lidar_calib_overlay}
\end{figure}
基于逐帧配准获得的外参估计序列，采用稳健筛选与拟合方法得到最终外参 $\hat{\mathbf{T}}_{l_2}^{l_1}$。同时，系统输出标定质量统计指标（包括有效帧比例与 fitness score 分布），为标定过程的可解释性与可靠性提供量化评估依据。图\ref{fig:lidar_calib_result}展示了自动标定过程的输出结果。

\begin{figure}[htbp]
  \centering
  \includegraphics[width=0.88\textwidth]{自动标定过程与结果输出.png}
  \caption{自动标定过程与结果输出}
  \label{fig:lidar_calib_result}
\end{figure}

本文最终采用的双雷达静态外参为
\begin{equation}
\label{eq:lidar_lidar_result}
\hat{\mathbf{T}}_{l_2}^{l_1}=
\begin{bmatrix}
-0.244146 & -0.00844596 & 0.969702 & 0.179165 \\
0.00174144 & -0.999964 & -0.00827110 & -0.00157140 \\
0.969737 & -0.00033068 & 0.244152 & -0.143879 \\
0 & 0 & 0 & 1
\end{bmatrix}
\end{equation}

式\eqref{eq:lidar_lidar_result}给出了从 $\{l_2\}$ 到 $\{l_1\}$ 的最终刚体变换。为确保该结果可直接用于后续在线融合，本文除检查标定前后点云叠加效果外，还在静止场景下统计重叠区域最近邻距离分布，并以最近邻距离均值不超过 $0.10\,\mathrm{m}$、95\% 分位不超过 $0.20\,\mathrm{m}$ 作为可用性判据。当该条件成立时，可认为双雷达在后续特征提取、配准和近场目标生成中具备可接受的几何一致性。

本章获得的 $\hat{\mathbf{T}}_{l_2}^{l_1}$ 将直接配置于 livox\_merge 的静态外参中，为第3章的双雷达实时融合提供输入。

\subsection{雷达-相机外参标定}\label{ux96f7ux8fbeux76f8ux673aux5916ux53c2ux6807ux5b9a}

雷达—相机外参标定旨在建立激光雷达点云与相机图像之间的精确投影关系，为第4章的雷视数据关联奠定几何基础。本文基于 Livox 开源标定工具\footnote{Livox\ Camera-LiDAR\ Calibration: \url{https://github.com/Livox-SDK/livox_camera_lidar_calibration}}，利用已知尺寸的立体靶标角点构建跨模态特征对应，通过非线性优化求解相机与雷达间的稳定刚体变换。

从数学形式上看，雷达—相机外参标定可表述为求解相机坐标系 $\{c\}$ 与基准雷达坐标系 $\{l_1\}$ 之间的刚体变换 $\mathbf{T}_{c}^{l_1} \in SE(3)$。本文采用统一的坐标变换记号：$\mathbf{T}_{b}^{a}$ 表示将坐标系 $\{b\}$ 下的点变换到坐标系 $\{a\}$。相机到雷达的变换矩阵定义为：
\begin{equation}
\label{eq:lidar_cam_extrinsic}
\mathbf{T}_{c}^{l_1}=
\begin{bmatrix}
\mathbf{R}_{c}^{l_1} & \mathbf{t}_{c}^{l_1}\\
\mathbf{0}^{\mathsf{T}} & 1
\end{bmatrix}\in SE(3)
\end{equation}
其中，$\mathbf{R}_{c}^{l_1} \in SO(3)$ 为旋转矩阵，$\mathbf{t}_{c}^{l_1} \in \mathbb{R}^3$ 为平移向量。

在实际投影计算中，需要将雷达点转换到相机坐标系。因此，运行时使用逆变换：
\begin{equation}
\label{eq:lidar_to_cam_tf_inverse}
\mathbf{T}_{l_1}^{c}=\left(\mathbf{T}_{c}^{l_1}\right)^{-1}=
\begin{bmatrix}
\mathbf{R}_{l_1}^{c} & \mathbf{t}_{l_1}^{c}\\
\mathbf{0}^{\mathsf{T}} & 1
\end{bmatrix}
\end{equation}

对于靶标角点在雷达坐标系中的三维坐标 $\mathbf{P}_{l_1}=[X_{l_1}\;Y_{l_1}\;Z_{l_1}]^{\mathsf{T}}$，其在相机坐标系下的坐标为：
\begin{equation}
\label{eq:lidar_to_cam}
\mathbf{P}_{c}=\mathbf{R}_{l_1}^{c}\,\mathbf{P}_{l_1}+\mathbf{t}_{l_1}^{c}=[X_c\;Y_c\;Z_c]^{\mathsf{T}}
\end{equation}

结合第2.2.1节标定的相机内参矩阵 $\mathbf{K}$ 与畸变模型，可预测角点在图像中的像素坐标：
\begin{equation}
\label{eq:lidar_cam_project}
\hat{u}=f_x\frac{X_c}{Z_c}+c_x,\qquad
\hat{v}=f_y\frac{Y_c}{Z_c}+c_y
\end{equation}
其中，$(f_x, f_y)$ 为等效焦距，$(c_x, c_y)$ 为主点坐标。设实际观测的像素坐标为 $(u,v)$，则单个角点的重投影残差为：
\begin{equation}
\label{eq:reproj_residual_uv}
\mathbf{r}=\begin{bmatrix}\hat{u}-u\\\hat{v}-v\end{bmatrix}
\end{equation}

基于多组靶标姿态（$k=1,\dots,K$）及各姿态的多个角点（$i=1,\dots,N_k$），构建非线性最小二乘优化问题：
\begin{equation}
\label{eq:cam_lidar_cost}
\displaystyle\min_{\mathbf{R}_{c}^{l_1},\mathbf{t}_{c}^{l_1}}\sum_{k=1}^{K}\sum_{i=1}^{N_k}\left\lVert \mathbf{r}_{k,i}\right\rVert_2^2
\end{equation}

求解最优外参 $\mathbf{T}_{c}^{l_1}$。为规避欧拉角奇异值问题，优化过程中采用四元数 $\mathbf{q}$ 参数化旋转，并施加单位约束。

标定流程采用“先数据录制、后离线处理”的工程化方式。首先，在开阔场地布置已知尺寸的矩形立体靶标，同步采集雷达点云（ROS bag）与相机图像，共采集12组不同距离（约3\,m）与姿态的数据，确保靶标在两种模态中均清晰可见。其次，在图像中提取靶标角点像素坐标 $(u,v)$（图\ref{fig:cam_lidar_corner_photo}），并在对应点云帧中提取相同角点的三维坐标 $\mathbf{P}_{l_1}$（图\ref{fig:cam_lidar_corner_picking}），采用一致的角点编号顺序（如顺时针）以避免特征对应错误，得到待标定的原始数据（图\ref{fig:cam_lidar_calib_dataset}）。接着，基于结构设计模型或安装测量提供外参初值 $\mathbf{T}_{c}^{l_1,(0)}$ 作为优化起点。然后固定已知内参，对外参进行非线性优化，根据重投影误差剔除异常数据后重新优化直至收敛。最后，将优化所得外参 $\mathbf{T}_{c}^{l_1}$ 用于点云到图像的投影，通过视觉对齐效果与点云着色验证标定精度（图\ref{fig:cam_lidar_calib_effect}）。
\begin{figure}[htbp]
  \centering
  \includegraphics[width=0.88\textwidth]{相机侧角点像素坐标获取.png}
  \caption{相机侧角点像素坐标获取}
  \label{fig:cam_lidar_corner_photo}
\end{figure}

\begin{figure}[H]
  \centering
  \includegraphics[width=0.88\textwidth]{雷达侧角点三维坐标获取.png}
  \caption{雷达侧角点三维坐标获取}
  \label{fig:cam_lidar_corner_picking}
\end{figure}

\begin{figure}[htbp]
  \centering
  \begin{minipage}{0.88\textwidth}
    \footnotesize
\begin{verbatim}
data
│
├── lidar/                              # LiDAR 原始数据（ROS bag）
│   ├── 0.bag
│   ├── 1.bag
│   ├── 2.bag
│   └── ...
│
├── pcdFiles/                           # 点云帧序列（PCD）
│   ├── 0.pcd                           # 按序号命名
│   ├── 1.pcd
│   ├── 2.pcd
│   └── ...
│
├── photo/                               # 图像帧序列（BMP）
│   ├── 0.bmp                            # 按序号命名
│   ├── 1.bmp
│   ├── 2.bmp
│   └── ...
│
├── parameters/                          # 参数文件
│   ├── intrinsic.txt                    # 相机内参
│   └── extrinsic.txt                    # 外参（或外参相关配置）
│
├── projection/                          # 投影/可视化输出
│   └── new.bmp
│
├── corner_lidar.txt                     # 角点/标定点（LiDAR）
└── corner_photo.txt                      # 角点/标定点（Image）
\end{verbatim}
  \end{minipage}
  \caption{雷达-相机外参标定数据准备}
  \label{fig:cam_lidar_calib_dataset}
\end{figure}

\begin{figure}[htbp]
  \centering
  \begin{subfigure}[t]{0.58\textwidth}
    \centering
    \includegraphics[height=6cm]{雷达点云投影到图像效果.png}
    \caption{雷达点云投影到图像效果}
    \label{fig:cam_lidar_projection}
  \end{subfigure}
  \hfill
  \begin{subfigure}[t]{0.4\textwidth}
    \centering
    \includegraphics[height=6cm]{雷达点云着色效果.png}
    \caption{雷达点云着色效果}
    \label{fig:cam_lidar_coloring}
  \end{subfigure}
  \caption{雷达-相机标定效果}
  \label{fig:cam_lidar_calib_effect}
\end{figure}

通过上述标定流程，最终求得相机到基准雷达1的静态外参矩阵为：
\begin{equation}
\label{eq:cam_lidar_result}
\mathbf{T}_{c}^{l_1}=
\begin{bmatrix}
0.785983 & -0.0208025 & -0.617898 & -0.0952186 \\
0.0338731 & 0.999382 & 0.0094418 & 0.117529 \\
0.617320 & -0.0283512 & 0.786201 & 0.0653585 \\
0 & 0 & 0 & 1
\end{bmatrix}
\end{equation}

该矩阵完整描述了相机坐标系 $\{c\}$ 相对于基准雷达坐标系 $\{l_1\}$ 的刚体安装关系。在实际应用中，通过逆变换 $\mathbf{T}_{l_1}^{c} = (\mathbf{T}_{c}^{l_1})^{-1}$ 可将雷达点云转换至相机坐标系，进而实现点云到图像的精确投影。这一变换关系为后续的雷视数据关联、跨模态融合与一致性验证提供了统一的几何基础。

\section{本章小结}\label{sec:ch2_summary}

本章面向吊钩端近场风险感知的工程部署需求，系统性地完成了多传感器系统的硬件集成、标定配准与数据接口规范，将“可部署”的传感器平台转化为“可计算”的统一观测空间。具体贡献包括：

（1）硬件架构与一体化集成。提出了双 LiDAR（Livox Mid-360）与工业相机（海康威视 MV-CS016-10UC）的异构配置方案，通过双雷达侧向对称布置（38°俯仰角）与相机垂直向下安装，实现了近场盲区互补与覆盖优化。基于 SolidWorks 完成了雷视融合一体化集成感知平台设计，将传感器、供电模块与边缘计算单元集中集成，为现场部署提供了结构可靠、可维护的硬件平台。

（2）多传感器外参标定与几何一致性建立。采用张正友标定法获取相机内参与畸变参数；提出了“粗配准初值构造+自动精配准优化”的两阶段双雷达外参标定方法，基于 GICP 点云配准与一致性筛选获得稳定外参；利用立体靶标角点构建跨模态特征对应，通过非线性优化求解雷达--相机外参。上述标定工作建立了点云与图像之间的精确投影关系，为后续雷视数据关联提供了统一的几何基础。

（3）通信链路与时统接口规范化。设计了“多源独立接入、边缘侧统一采集、系统总线解耦发布”的通信链路架构，通过图像压缩发布与队列长度控制降低链路时延与抖动。采用 PTPv2 软件时间戳建立统一时基，结合融合层时间对齐、延迟补偿与时间窗门控策略，实现了跨模态观测的软件时间一致性。

（4）分层式空间参考体系与 TF 接口规范。提出了“动态定位主链”与“静态感知支链”相分离的坐标系架构，基于 ROS TF 框架制定了统一的坐标系定义、变换链查询与外参管理规范，确保多源观测在同一几何框架下的可解释性与可复现性。

本章工作为第3章的融合点云定位建图与危险区域检测算法提供了稳定、统一且可维护的数据输入基础，同时也为第4章的雷视关联与风险判定奠定了时空一致的跨模态融合前提。

